\section{Two certificate mechanisms}
\label{rev:certificate-mechanisms}

The two constructions below impose arithmetic restrictions through
explicit auxiliary identities. Pad\'e approximation separates two
rationally related values; rational telescoping classifies identities
within two fixed deformation families. Neither construction is known
to be necessary for an arbitrary standard identity.

\subsection{An explicit exclusion for two rationally proportional multipliers}

The preceding specialization arguments exclude nearby pairs
asymptotically.  A direct Pad\'e construction gives an explicit threshold
when the ratio of the two multipliers is rational.  We use the polynomial
construction of David--Hirata-Kohno--Kawashima
\cite[Proposition~3.6]{DHK2022}, and prove the determinant and all
growth estimates needed here.

\begin{theorem}[Explicit two-point independence]
\label{thm:C-pade-explicit}
For the four standard levels, define
\[
 (\Lambda_1,\Lambda_2,\Lambda_3,\Lambda_4)
   =\left(\frac{46}{5},\frac{22}{3},\frac{13}{2},6\right),
 \qquad U=6\log2+3(7\log7-6\log6),
\]
and put $T_\ell=U+6\Lambda_\ell+6\log R_\ell$.
Let $a_1,a_2$ be distinct nonzero integers and
$M=\max(|a_1|,|a_2|)$.  If the integer $b$ satisfies
\begin{equation}\label{eq:C-pade-threshold}
 b>M^7\exp(T_\ell),
\end{equation}
then the following seven numbers are linearly independent over $\Q$:
\[
 1,\quad F_s(a_i/b),\quad
 (\theta F_s)(a_i/b),\quad(\theta^2F_s)(a_i/b)
 \qquad(i=1,2),\qquad s=s_\ell.
\]
In particular, this concerns ordinary absolutely convergent series.
\end{theorem}

We first establish the required determinant independently.

\begin{lemma}[The two-point Pad\'e determinant]
\label{lem:C-pade-determinant}
Fix $s\in\{1/6,1/4,1/3,1/2\}$ and put
\[
 d_j=\frac{(1/2)_j(s)_j(1-s)_j}{(j!)^3},\qquad c_k=d_{k+1}.
\]
Let $a_1,a_2$ be distinct nonzero rational numbers, $n\ge1$, and
$0\le h\le6$.  Write
\[
 G_h(t)=t^h(t-a_1)^{3n}(t-a_2)^{3n}
       =\sum_k g_{h,k}t^k,
 \qquad
 P_h(t)=\sum_k\frac{g_{h,k}}{c_k}\binom{k+n}{n-1}^{\!3}t^k.
\]
For $q=0,1,2$, define the rational linear functional
\[
 \psi_{i,q}(t^k)=(k+1)^q c_k a_i^{k+1},\qquad
 Q_{h,i,q}(z)=\psi_{i,q}\!\left(
           \frac{P_h(z)-P_h(t)}{z-t}\right).
\]
Then the seven columns $(P_h(z),Q_{h,i,q}(z))$, $0\le h\le6$,
form an invertible matrix for every $z$.  Moreover, if
$\mathcal F_0=F_s-1$, $\mathcal F_1=\theta F_s$, and
$\mathcal F_2=\theta^2F_s$, then
\begin{equation}\label{eq:C-pade-remainder}
 P_h(z)\mathcal F_q(a_i/z)-Q_{h,i,q}(z)
   =\sum_{j\ge0}\frac{\psi_{i,q}(t^{n+j}P_h(t))}{z^{n+j+1}}.
\end{equation}
The equality is formal at infinity and convergent for
$|z|>\max(|a_1|,|a_2|)$.
\end{lemma}

\begin{proof}
Write $\theta=t\,d/dt$ and
\[
 A(X)=(X+3/2)(X+1+s)(X+2-s),\qquad B(X)=(X+1)^3.
\]
The recurrence $c_{k+1}/c_k=A(k)/B(k+1)$ gives the construction in
\cite[Proposition~3.6]{DHK2022}.  We verify the features needed here.
Expanding the divided difference gives the right side of
\eqref{eq:C-pade-remainder} initially with $t^{j}$ and $z^{j+1}$,
starting at $j=0$.  For $0\le j<n$, the recurrence and
$t^jH(\theta)=H(\theta-j)t^j$ give the exact identity
\[
 \psi_{i,q}(t^jP_h)=\frac{a_i}{(n-1)!^3}
 \operatorname{ev}_{a_i}(\theta+1)^q
 \prod_{r=1}^{j}A(\theta-r)
 \prod_{r=1}^{n-1-j}B(\theta+r)\,[t^jG_h(t)].
\]
The differential operator has order at most
\[
 q+3j+3(n-1-j)\le3n-1
\]
on $t^{h+j}(t-a_1)^{3n}(t-a_2)^{3n}$.  It vanishes, proving the
formal equality.  Since $c_k$ and $c_k^{-1}$ have polynomial growth,
the stated convergence follows.

Order the lower rows by $(i,q)=(1,2),(1,1),(1,0),(2,2),(2,1),(2,0)$
and denote the determinant by $\Delta(z)$.  Replacing these six rows
by their remainders does not change the determinant.  Each remainder
has order at least $n+1$ at infinity, whereas $\deg P_h=6n+h$.
Thus $\Delta(z)$ is constant.  Only the leading term in column
$h=6$ can contribute, so
\begin{equation}\label{eq:C-pade-delta-moments}
 \Delta=L_n\det\bigl(\psi_{i,q}(t^nP_h(t))\bigr)_{
                    (i,q),\,0\le h\le5},\qquad
 L_n=c_{6n+6}^{-1}\binom{7n+6}{n-1}^{\!3}.
\end{equation}

Put $H(X)=\prod_{k=1}^{n}A(X-k)$ and
$V_h=t^{n+h}(t-a_1)^{3n}(t-a_2)^{3n}$.  The same recurrence,
coefficient by coefficient, gives
\[
 \psi_{i,q}(t^nP_h(t))
  =\frac{a_i}{(n-1)!^3}
    \operatorname{ev}_{a_i}H(\theta)(\theta+1)^{q-3}V_h.
\]
Expand $H(X)$ in powers of $X+1$.  Terms with a nonnegative
power of $\theta+1$ vanish at $a_i$, since their differential
order is at most $3n-1$.  The remaining three rows are a triangular
combination, with diagonal coefficient $H(-1)$, of the functionals
\[
 \varphi_{a,r}(t^k)=\frac{a^k}{(k+1)^{r+1}},\quad r=0,1,2,
 \qquad
 H(-1)=\prod_{k=1}^n(1/2-k)(s-k)(1-s-k)\ne0.
\]
For integers $u\ge0$, define
\[
 C_{n,u}(a,b)=\det\left(
  \varphi_{a_i,r}\bigl(t^{u+h}(t-a)^{3n}(t-b)^{3n}\bigr)
                         \right)_{(i,r),\,0\le h\le5},
 \quad (a_1,a_2)=(a,b).
\]
It follows that the moment determinant in
\eqref{eq:C-pade-delta-moments} equals
\begin{equation}\label{eq:C-pade-moment-reduction}
 \frac{(a_1a_2)^3H(-1)^6}{(n-1)!^{18}}C_{n,n}(a_1,a_2).
\end{equation}

Here is a direct factorization of $C_{n,u}$.  In the ordered variables
$(X_0,X_1,X_2,Y_0,Y_1,Y_2)$, let
\[
 \mathcal P=\prod_Z Z^u(Z-a)^{3n}(Z-b)^{3n}
               \prod_{i<j}(Z_j-Z_i).
\]
Expanding the determinant expresses $C_{n,u}$ as the product of the
six functionals $\varphi_{a,r}^{X_r},\varphi_{b,r}^{Y_r}$ applied
to $\mathcal P$.  Its degree is $6u+36n+15$.
Scaling the three $X$ variables by $a$ shows divisibility by
$a^{3u+9n+3}$, and scaling the $Y$ variables gives the same power
of $b$.

We next prove divisibility by $(b-a)^{18n+9}$.  Away from $a=0$,
\[
 \partial_a\varphi_{a,r}(f)
 =\varphi_{a,r}(\partial_af)
    +a^{-1}\varphi_{a,r}(\theta f).
\]
After $v$ derivatives with respect to $a$, each term therefore applies
the six moment functionals to
\[
 \prod_{r=0}^2\theta_{X_r}^{v_r}\partial_a^j\mathcal P,
 \qquad j+v_0+v_1+v_2\le v,
\]
up to a scalar regular at any fixed nonzero $a$.  Expand each power
of $\theta_{X_r}$ in the monic polynomial basis
\[
 T_{r,k}(\theta)=
 \begin{cases}
  (\theta+1)^k,&k\le r+1,\\
  (\theta+1)^{r+1}\theta^{k-r-1},&k\ge r+1.
 \end{cases}
\]
Only $k_r\le v_r$ occur.  On setting $b=a$, if some
$k_r<r+1$, the composed $X_r$ functional equals the
$Y_{r-k_r}$ functional.  The polynomial $\partial_a^j\mathcal P$
is antisymmetric in these variables: the factor outside its
Vandermonde is symmetric, even after parameter differentiation.
The two identical functionals consequently give zero.

If all $k_r\ge r+1$, cancellation of the three inverse Euler
operators leaves evaluation at $X_0=X_1=X_2=a$ after an operator
of total differential order $q=\sum_r k_r-6$.
Before parameter differentiation and on $b=a$, the polynomial
has total vanishing order at least $18n+3$ in
$(X_0-a,X_1-a,X_2-a)$: the three powers contribute $18n$ and
their Vandermonde contributes $3$.  The $j$ parameter derivatives
lower that order by at most $j$.  Thus the evaluation is zero when
\[
 q+j\le v-6<18n+3.
\]
All derivatives of order $v<18n+9$ vanish on the diagonal, proving
the asserted factor.  Comparing degrees gives
\begin{equation}\label{eq:C-pade-coalescence}
 C_{n,u}(a,b)=\kappa_{n,u}
       (ab)^{3u+9n+3}(b-a)^{18n+9}.
\end{equation}

It remains to show $\kappa_{n,u}\ne0$.  Put
\[
 J_{n,u}=\det\left(
  \int_0^1 t^{u+h}(t-1)^{3n}
          \frac{(-\log t)^r}{r!}\,dt
                  \right)_{0\le r,h\le2}.
\]
The determinant expansion and symmetrization express this as $1/3!$
times an integral over $(0,1)^3$ whose integrand is
\[
 \prod_{j=1}^3t_j^u(t_j-1)^{3n}
   \det(t_j^h)_{h,j}
   \det\!\left(\frac{(-\log t_j)^r}{r!}\right)_{r,j}.
\]
On $0<t_1<t_2<t_3<1$, the ordinary Vandermonde is positive and
the logarithmic Vandermonde has sign $(-1)^3$.  The weight also
has constant sign.  Their product is symmetric and nonzero off
the diagonals, so the integral is nonzero.  In
\eqref{eq:C-pade-coalescence}, set $b=1$, divide by
$a^{3u+9n+3}$, and let $a\to0$ after scaling the $X$ variables.
The two three-variable blocks separate, giving
\[
 \kappa_{n,u}=(-1)^{9n}J_{n,u}J_{n,u+3n+3}\ne0.
\]
Equations \eqref{eq:C-pade-delta-moments}--
\eqref{eq:C-pade-coalescence} now give
\[
 \Delta=\frac{L_nH(-1)^6\kappa_{n,n}}{(n-1)!^{18}}
       (a_1a_2)^{12n+6}(a_2-a_1)^{18n+9}\ne0,
\]
as required.
\end{proof}

\begin{proof}[Proof of Theorem~\ref{thm:C-pade-explicit}]
We first bound the common denominators in the preceding construction.
Write $d_j=Q_j/R^j$, where the four integer sequences $Q_j$ are
\[
 \frac{(6j)!}{(3j)!(j!)^3},\qquad
 \frac{(4j)!}{(j!)^4},\qquad
 \frac{(3j)!(2j)!}{(j!)^5},\qquad
 \frac{(2j)!^3}{(j!)^6}.
\]
Put $\mathcal L(x)=\operatorname{lcm}(1,\ldots,\lfloor x\rfloor)$.
Legendre's factorial formula expresses the valuation at each prime $p$
as a sum, over $p^v$, of the step functions
\[
 \lfloor6u\rfloor-\lfloor3u\rfloor,\quad
 \lfloor4u\rfloor,\quad
 \lfloor3u\rfloor+\lfloor2u\rfloor,\quad
 3\lfloor2u\rfloor,\qquad u=\{j/p^v\}.
\]
The thresholds for contributions $1,2,3$ are respectively
\[
 (1/6,1/2,5/6),\quad(1/4,1/2,3/4),\quad
 (1/3,1/2,2/3),\quad(1/2,1/2,1/2).
\]
For $j\le N$, a contribution of at least $r$ implies
$p^v\le N/t_r$, with the corresponding threshold $t_r$.
Consequently every $Q_j$, $1\le j\le N$, divides the indicated
integer in this list:
\[
 \mathcal L(6N)\mathcal L(2N)\mathcal L(6N/5),\quad
 \mathcal L(4N)\mathcal L(2N)\mathcal L(4N/3),
\]
\[
 \mathcal L(3N)\mathcal L(2N)\mathcal L(3N/2),\qquad
 \mathcal L(2N)^3.
\]
Let $D_N$ and $E_N$ be the least common denominators of, respectively,
$\{d_j^{-1}:1\le j\le N\}$ and $\{d_j:1\le j\le N\}$.
Since $E_N\mid R^N$, the prime number theorem gives
\begin{equation}\label{eq:C-pade-denominators}
 \log D_N\le\Lambda_\ell N+o(N),\qquad
 \log E_N\le N\log R_\ell.
\end{equation}

Take the polynomials from Lemma~\ref{lem:C-pade-determinant} and
$N=6n+7$.  The binomial factors and $a_i$ are integers; $D_N$
therefore clears the coefficients of every $P_h$.  Applying
$\psi_{i,q}$ to the divided differences introduces only coefficients
$c_k$ with $k+1\le N$, so $D_NE_N$ clears all the $Q_{h,i,q}$
as well.  Evaluation at integer $b$ introduces no further denominator.
By \eqref{eq:C-pade-denominators},
\begin{equation}\label{eq:C-pade-cost}
 \log(D_NE_N)\le6n(\Lambda_\ell+\log R_\ell)+o(n).
\end{equation}

We estimate the remainders directly.  The coefficient of $t^k$ in
$G_h(t)$ has absolute value at most
$2^{6n}M^{6n+h-k}$.  Its diagonal binomial factor is at most
$\binom{7n+6}{n-1}^{3}$.  Stirling's formula and the polynomial
growth of $c_k,c_k^{-1}$ give, for a fixed integer $K$, a bound
\[
 |\psi_{i,q}(t^{n+j}P_h(t))|
 \le e^{Un+o(n)}(n+j+1)^K M^{7n+h+j+1}.
\]
Indeed, the exponent $-k$ from $G_h$ cancels the $k$ in
$a_i^{k+n+j+1}$; the remaining coefficient dependence is polynomial.
As $b>M$, summing \eqref{eq:C-pade-remainder} gives, uniformly
over the seven choices of $h$,
\begin{equation}\label{eq:C-pade-error}
 |P_h(b)\mathcal F_q(a_i/b)-Q_{h,i,q}(b)|
 \le\exp\{n[U+7\log M-\log b]+o(n)\}.
\end{equation}
The polynomial in $j$ is harmless in this fixed-ratio geometric sum.

Suppose a nonzero rational relation exists among
$1,\mathcal F_q(a_i/b)$.  Clearing denominators gives integer
coefficients $\lambda_0,\lambda_{i,q}$.  By the nonzero determinant,
for every $n$ at least one $h$ gives a nonzero rational number
\[
 L_h=\lambda_0P_h(b)+\sum_{i,q}\lambda_{i,q}Q_{h,i,q}(b).
\]
The assumed relation writes $L_h$ as minus the corresponding sum
of remainders.  Since $D_NE_NL_h$ is a nonzero integer,
\eqref{eq:C-pade-cost} and \eqref{eq:C-pade-error} imply
\[
 1\le|D_NE_NL_h|
 \le\exp\{n[T_\ell+7\log M-\log b]+o(n)\}\longrightarrow0,
\]
contradicting \eqref{eq:C-pade-threshold}.  The replacement
$\mathcal F_0=F_s-1$ is an invertible rational change of the value
vector and proves the stated result.
\end{proof}

\begin{corollary}[An explicit opposite-sign exclusion]
\label{cor:C-pade-opposite}
At level $\ell$, let $b$ be a positive integer.  If
\[
 b\ge10^{50},\quad10^{40},\quad10^{35},\quad10^{33}
 \qquad\text{for }\ell=1,2,3,4\text{ respectively},
\]
then the two denominators $C=R_\ell b$ and $C=-R_\ell b$ cannot
both support standard identities with nonzero algebraic multipliers
whose ratio is rational.  In particular, they cannot both support
identities with rational nonzero multipliers.
\end{corollary}

\begin{proof}
Take $(a_1,a_2)=(1,-1)$, so $M=1$.  The four displayed integers
exceed $\exp(T_\ell)$.  For an exact check, use $e<11/4$ and
\[
 \exp(T_\ell)
 =\frac{2^6 7^{21}R_\ell^6}{6^{18}}e^{6\Lambda_\ell}
 <\frac{2^6 7^{21}R_\ell^6}{6^{18}}
       (11/4)^{\lceil6\Lambda_\ell\rceil};
\]
comparison with each indicated power of ten is an integer inequality.
Subtracting a rational multiple of one asserted identity from the
other would give a nonzero rational relation among the two copies
of $F_s,\theta F_s$, contrary to
Theorem~\ref{thm:C-pade-explicit}.
\end{proof}

This result imposes no degree or height bound on the individual
multipliers; it restricts their ratio to $\Q$.  It excludes a pair
of identities and does not exclude an isolated non-CM identity.
The seven-number independence proved here is over $\Q$, not over
$\Qbar$.  The determinant argument therefore does not supply
unconditional CM forcing.

\subsection{Complete converses within two rational-certificate families}
\label{sec:C-wz-converses}

A rational Wilf--Zeilberger certificate supplies a second way to impose
arithmetic restrictions without assuming a period conjecture.  The
restriction is to an explicitly specified deformation of the summand.
We do not assume, or prove, that an arbitrary scalar identity admits
such a deformation.

We first allow three affine parameter shifts and a varying linear
weight.  For $0<s<1$, $z\in\C\setminus\{0,1\}$, and
$\lambda\in\C^\times$, put
\begin{align}
 U(n,k)&=z^n\frac{(1/2+k)_n(s)_n(1-s)_n}
 {(1+k)_n(1+2k)_n(1)_n},\label{eq:C-wz-affine-U}\\
 H(k)&=\lambda^k
 \frac{\Gamma(k+1)^2\Gamma(s)\Gamma(1-s)}
 {\Gamma(k+s)\Gamma(k+1-s)},\qquad T(n,k)=\frac{U(n,k)}{H(k)}.
 \label{eq:C-wz-affine-H}
\end{align}
Here a determination of $\lambda^k$ is fixed; only its unit-shift
ratio is used in the rational calculation.  At $k=0$ one has
$H(0)=1$ and $U(n,0)=c_s(n)z^n$.  The level-three specialization of
this deformation occurs in Guillera's parameterized WZ
identity.\footnote{J.~Guillera, \emph{On WZ-pairs which prove
Ramanujan series}, arXiv:0904.0406v3, equation~(37),
\url{https://arxiv.org/abs/0904.0406v3}.  That formula is a forward
identity; the rational-certificate converse below is proved directly.}
For $P(n,k)=A+Bn+Dk$, a rational certificate means an arbitrary
$R\in\C(n,k)$ such that
\begin{equation}\label{eq:C-wz-affine-identity}
\begin{split}
 P(n,k+1)T(n,k+1)-P(n,k)T(n,k)
 ={}&R(n+1,k)T(n+1,k)\\
 &-R(n,k)T(n,k).
\end{split}
\end{equation}
The equation is interpreted after division by $T(n,k)$ as a rational
identity.  No bound on the numerator or denominator degrees of $R$
is imposed.

\begin{theorem}[An affine-shift certificate converse]
\label{thm:C-wz-affine}
Assume $(A,B)\ne(0,0)$.  A rational certificate
\eqref{eq:C-wz-affine-identity} exists precisely in the following two
cases, up to a common nonzero scalar multiplying $A,B,D,R$:
\[
\begin{array}{c|c|c|c}
 s&z&\lambda&(A,B,D)\\\hline
 1/3\text{ or }2/3&1/2&1&(1/6,1,1)\\
 1/3\text{ or }2/3&-4&1&(4/15,1,2/5).
\end{array}
\]
The corresponding certificates, respectively, are unique and equal to
\[
 R=\frac{2n(3k+3n+2)}{3(2k+n+1)},\qquad
 R=\frac{n(12k+3n+5)}{15(2k+n+1)}.
\]
Among the four standard families, ordinary convergence leaves only
$z=1/2$, $s=1/3$, $C=216$, with primitive coefficient pair $(1,6)$.
Its elliptic fiber has CM of discriminant $-24$.
\end{theorem}
\begin{proof}
Set $L=4/\lambda$.  Direct unit shifts give
\begin{align*}
 q:=\frac{T(n+1,k)}{T(n,k)}
 &=\frac{z(n+k+1/2)(n+s)(n+1-s)}
 {(n+1)(n+k+1)(n+2k+1)},\\
 K:=\frac{T(n,k+1)}{T(n,k)}
 &=\frac{L(n+k+1/2)(k+s)(k+1-s)}
 {(n+k+1)(n+2k+1)(n+2k+2)}.
\end{align*}
Make the invertible rational substitution $R=nX/(n+2k+1)$.
Clearing denominators yields
\begin{equation}\label{eq:C-wz-affine-polynomial}
\begin{split}
 &z(n+k+1/2)(n+s)(n+1-s)X(n+1,k)\\
 &\quad-n(n+k+1)(n+2k+2)X(n,k)\\
 &=L(n+k+1/2)(k+s)(k+1-s)(A+Bn+D(k+1))\\
 &\quad-(A+Bn+Dk)(n+k+1)(n+2k+1)(n+2k+2).
\end{split}
\end{equation}
We show that $X$ is a polynomial in $n$.  Regard $k$ as
transcendental and work over an algebraic closure of $\C(k)$.
If $X$ has a pole, choose the first and last poles in one finite
chain under integer translation.  Cancellation at the endpoints
requires a root of the first coefficient of $X$ in
\eqref{eq:C-wz-affine-polynomial}, shifted by a positive integer,
to equal a root of the second coefficient.  Those root sets are
\[
 \{-k-1/2,-s,s-1\},\qquad \{0,-k-1,-2k-2\}.
\]
The constant roots in the first set lie strictly between $-1$ and
$0$.  The only nonconstant pair having equal coefficients of $k$
has difference $-1/2$.  All remaining comparisons have unequal
coefficients of the transcendental $k$.  Thus none gives the required
positive integral difference, and $X$ has no pole.

If $d=\deg_nX$, the left side of
\eqref{eq:C-wz-affine-polynomial} has degree $d+3$, with leading
coefficient $(z-1)\operatorname{lc}_n(X)$; its right side has degree
at most four.  Hence $X=u(k)n+v(k)$.  The coefficients of $n^4,n^3$
force
\begin{equation}\label{eq:C-wz-affine-uv}
\begin{split}
 u&=-\frac B{z-1},\\
 v&=-\frac{2A(z-1)+8Bkz-4Bk+3Bz-2B+2Dk(z-1)}
 {2(z-1)^2}.
\end{split}
\end{equation}
Let $E(n,k)$ be $(z-1)^2$ times the left side minus the right side
of \eqref{eq:C-wz-affine-polynomial}, after
\eqref{eq:C-wz-affine-uv} is substituted, and write
$E_{ij}=[n^ik^j]E$.  The classification is now a finite coefficient
calculation, as follows.

First suppose $B\ne0$ and divide by $B$ so that $B=1$.  Two
coefficients are
\[
 E_{04}=-D(L-4)(z-1)^2,\qquad
 E_{03}=-\tfrac12(2A+5D)(L-4)(z-1)^2.
\]
If $L\ne4$, they force $A=D=0$.  At this specialization
$E_{12}-3E_{22}/2=z$, contradicting $z\ne0$.  Thus $L=4$ and
$\lambda=1$.  Next,
\[
 E_{22}=2(2z-1)((z-1)D+2).
\]
If $z=1/2$, the equation $E_{21}=0$ gives $D=1$.  With
$w=s^2-s$, the coefficients
\[
 E_{20}=(-2A+12w+3)/16,\qquad E_{02}=(4A+3w)/4
\]
give $w=-2/9$ and $A=1/6$, hence $s=1/3$ or $2/3$.

If $z\ne1/2$, one instead has $D=-2/(z-1)$.  Then
\begin{align*}
 E_{21}&=\frac{4A(z-1)(2z-1)-z^2+7z-4}{2},\\
 E_{12}&=2A(z-1)(2z-1)+5z-4.
\end{align*}
Their difference is $(z+4)(z-1)/2$.  Consequently $z=-4$, and
$E_{12}=0$ gives $A=4/15$.  Finally,
\[
 E_{02}=4(z-1)\{A(z-1)+(s^2-s)(z-2)\}
\]
gives $s^2-s=-2/9$ once more.

If $B=0$, then $A\ne0$.  The same $E_{04},E_{03}$ force $L=4$.
In this case
\begin{align*}
 E_{22}&=2D(z-1)(2z-1),\\
 E_{21}&=\tfrac12\{4A(z-1)(2z-1)+D(z-1)(5z-2)\}.
\end{align*}
Their vanishing, with $z\ne1$, implies $D=0,z=1/2$.  But then
$E_{20}=-A/8\ne0$.  Thus no additional case occurs.

Substitution in the full polynomial identity proves sufficiency and
gives the two displayed certificates.  For uniqueness, subtract two
certificates.  The same pole argument makes their difference
correspond to a polynomial $X$, whose nonzero leading term could not
vanish in the homogeneous version of
\eqref{eq:C-wz-affine-polynomial} because $z\ne1$.

At $k=0$, the $z=-4$ survivor has $B\ne0$ and its summands have
magnitude proportional to $4^n n^{-1/2}$; they do not tend to zero.
It cannot define an ordinarily convergent standard identity.  The
other survivor has $C=108/(1/2)=216$.  Formula~\eqref{eq:P3} gives
\[
 P_3(J,216)=J^2-4834944J+14670139392=H_{-24}(J).
\]
Both corresponding elliptic invariants are therefore singular moduli
of discriminant $-24$.
\end{proof}

For completeness the convergent identity follows directly from the
certificate, without analytic continuation in its auxiliary
parameter.  This also specifies the analytic step that a rational
certificate alone does not supply.
\begin{corollary}\label{cor:C-wz-affine-value}
The ordinary sum at the convergent survivor is
\[
 (1+6\thetaop)F_{1/3}(1/2)=\frac{3\sqrt3}{\pi}.
\]
\end{corollary}
\begin{proof}
Use $(A,B,D)=(1,6,6)$ and $s=1/3,z=1/2,\lambda=1$.
For every nonnegative integer $k$, the series
\[
 S(k)=\sum_{n\ge0}(1+6n+6k)\frac{U(n,k)}{H(k)}
\]
converges absolutely.  The certificate has zero boundary at $n=0$
and at infinity: its rational factor has polynomial growth, whereas
the summand ratio tends to $1/2$.  Telescoping therefore gives
$S(k+1)=S(k)$ on the nonnegative integers.

For integers $k\ge1$ and $0<s<1$,
\[
 \frac{(k+1/2)_n}{(k+1)_n}\le1,\qquad
 \frac{(s)_n}{n!}\le1,\qquad
 \frac{(1-s)_n}{(1+2k)_n}\le1.
\]
Thus $0<U(n,k)\le2^{-n}$; for fixed $n\ge1$ one also has
$U(n,k)\to0$ as $k\to\infty$.  The terms
$(1+6n+6k)U(n,k)/k$ are dominated by $(7+6n)2^{-n}$, a summable
bound independent of the integer $k\ge1$.  Dominated convergence
and the gamma-ratio asymptotic give, respectively,
\begin{align*}
 \frac1k\sum_{n\ge0}(1+6n+6k)U(n,k)&\longrightarrow6,\\
 \frac{H(k)}k&\longrightarrow
 \Gamma(s)\Gamma(1-s)=\frac\pi{\sin\pi s}.
\end{align*}
Hence $S(0)=\lim_{k\to\infty}S(k)=6\sin(\pi/3)/\pi$.
Only integer values of $k$ have been used; no claim that a
one-periodic analytic function must be constant enters the proof.
\end{proof}

The second converse allows an unrestricted rational dependence of
the auxiliary coefficients and of the normalization ratio.  It
shows that such a change does not enlarge the earlier terminating
mechanism.
\begin{theorem}[Rational auxiliary gauges of the terminating family]
\label{thm:C-wz-EZ-gauge}
Let $0<s<1$, let $(a,b,t)$ be a permutation of $(1/2,s,1-s)$,
and let $z\in\C\setminus\{0,1\}$.  Set
\[
 U(n,k)=z^n\frac{(a)_n(b)_n(-k)_n}
 {(1)_n^2(1+t+k)_n},\qquad T(n,k)=U(n,k)/H(k),
\]
where $h(k)=H(k+1)/H(k)\in\C(k)^\times$.  Let
$A,B\in\C(k)$ be not both identically zero.  An arbitrary rational
$R(n,k)$ satisfying
\begin{equation}\label{eq:C-wz-EZ-gauge-identity}
\begin{split}
 &[A(k+1)+nB(k+1)]T(n,k+1)-[A(k)+nB(k)]T(n,k)\\
 &\hspace{18mm}=R(n+1,k)T(n+1,k)-R(n,k)T(n,k)
\end{split}
\end{equation}
exists if and only if
\begin{gather*}
 s=a=b=t=1/2,\qquad z=-1,\qquad A(k)=B(k)/4,\\
 h(k)=\frac{k+3/2}{k+1}\frac{B(k+1)}{B(k)},\qquad
 R(n,k)=-\frac{B(k)n^2}{2(k+1-n)},
\end{gather*}
with $B(k)\ne0$ as a rational function.  The certificate is unique.
\end{theorem}
\begin{proof}
Put $x=k+1$ and $L(k)=(k+1)(k+1+t)/h(k)$.  The shift ratios are
\begin{align*}
 \frac{T(n+1,k)}{T(n,k)}&=
 \frac{z(n+a)(n+b)(n-k)}{(n+1)^2(n+k+1+t)},\\
 \frac{T(n,k+1)}{T(n,k)}&=
 \frac{L(k)}{(k+1-n)(n+k+1+t)}.
\end{align*}
Substitute $R=n^2X/(k+1-n)$ and clear denominators to obtain
\begin{equation}\label{eq:C-wz-EZ-polynomial}
\begin{split}
 &-z(n+a)(n+b)(k+1-n)X(n+1,k)\\
 &\quad-n^2(n+k+1+t)X(n,k)\\
 &=L(k)\{A(k+1)+nB(k+1)\}\\
 &\quad-\{A(k)+nB(k)\}(k+1-n)(n+k+1+t).
\end{split}
\end{equation}
The same pole-endpoint argument applies, now to the root sets
$\{-a,-b,k+1\}$ and $\{0,-k-1-t\}$.  No root of the first set
reaches a root of the second by a positive integer translation:
$0<a,b<1$, and the nonconstant comparisons have incompatible
coefficients of $k$.  Thus $X$ is polynomial in $n$.  The leading
coefficient on the left is $(z-1)\operatorname{lc}_n(X)$ and the
right side has degree at most three, so $X$ is constant in $n$.
If $B=0$ identically, comparison of degrees forces $X=0$, and the
$n^2$ coefficient then forces $A=0$, a contradiction.  Consequently
\[
 X=\frac{B(k)}{z-1},\qquad
 \frac{A(k)}{B(k)}=w(x):=
 \frac{z(a+b-t)-(z+1)x}{z-1}.
\]
The coefficients of $n$ and $1$ in
\eqref{eq:C-wz-EZ-polynomial} give
\[
 L(k)\frac{B(k+1)}{B(k)}=M(x),\qquad
 L(k)\frac{A(k+1)}{B(k)}=N(x),
\]
where
\begin{align*}
 M(x)&=\frac{z(ab-(a+b)x)}{z-1}-tw(x)+x(x+t),\\
 N(x)&=w(x)x(x+t)-\frac{zabx}{z-1}.
\end{align*}
Their necessary compatibility condition is the polynomial identity
\[
 \mathcal Q(x):=(z-1)^2\{N(x)-w(x+1)M(x)\}=0.
\]
If $t=1/2$, take $a=s,b=1-s$.  The coefficients of $x^2,x,1$
in $\mathcal Q$ are, respectively,
\[
 \frac{z^2-1}{2},\qquad
 -\frac{z(2s-1)^2}{2},\qquad
 -\frac{z(2s-1)^2(z+2)}8.
\]
They imply $s=1/2,z=-1$.  Otherwise, symmetry in $a,b$ permits
$a=1/2,b=1-t$.  The coefficients are
\[
 \frac{(z+1)(4tz+2t-z-2)}2,
\]
\[
 \frac{z(2t-1)(12tz+8t-3z+2)}4,
 \qquad
 \frac{z(2t-1)^2(4tz-z+2)}4.
\]
If $t\ne1/2$, the last gives $4tz-z=-2$ and the middle then gives
$8t-4=0$, a contradiction.  The sole survivor is therefore
$s=a=b=t=1/2,z=-1$.  At this survivor $w=1/4$ and $M=x^2$,
which give the stated formulas for $A,h,R$.  Their direct
substitution proves sufficiency.  The homogeneous pole and leading
coefficient argument proves uniqueness as before.
\end{proof}

The last theorem identifies all its certificates with rational
auxiliary gauges of Bauer's mechanism.  A one-periodic factor in
$H$ is invisible to its rational unit-shift ratio, and is not fixed
by the algebraic classification.  An analytic specialization still
requires regularity and the relevant convergence argument.  Neither
of these two theorems supplies a converse from a bare scalar value
to the existence of a rational WZ certificate.  Their CM conclusions
are complete within the displayed certificate families, while
unconditional CM forcing for an arbitrary standard identity remains
unresolved.

